\documentclass[11pt]{article}

\usepackage[margin=1in]{geometry}
\usepackage{amsmath}
\usepackage{amssymb}
\usepackage{booktabs}
\usepackage[hidelinks]{hyperref}

\newcommand{\softmax}{\operatorname{softmax}}

\title{\textbf{Exercise 01: Scaled Dot-Product Self-Attention}}
\author{}
\date{}

\begin{document}
\maketitle

\begin{center}
\emph{A complete numerical walkthrough of one single-head self-attention pass.}
\end{center}

\section*{Setup}

Using row-token notation,
\[
X=
\begin{bmatrix}
1&0\\
0&1\\
1&1
\end{bmatrix},
\qquad
W_Q=
\begin{bmatrix}1&0\\0&1\end{bmatrix},
\]
\[
W_K=
\begin{bmatrix}1&1\\0&1\end{bmatrix},
\qquad
W_V=
\begin{bmatrix}1&0\\1&1\end{bmatrix},
\qquad d_k=2.
\]

The attention operation is
\[
\mathrm{Attention}(Q,K,V)
=\softmax\left(\frac{QK^\top}{\sqrt{d_k}}\right)V.
\]

\section*{1. Compute $Q$, $K$, and $V$}

Because $W_Q=I_2$,
\[
Q=XW_Q=
\begin{bmatrix}
1&0\\
0&1\\
1&1
\end{bmatrix}.
\]

The key projection is
\[
K=XW_K=
\begin{bmatrix}
1&1\\
0&1\\
1&2
\end{bmatrix}.
\]

The value projection is
\[
V=XW_V=
\begin{bmatrix}
1&0\\
1&1\\
2&1
\end{bmatrix}.
\]
All three matrices have shape $3\times2$.

\section*{2. Compute the score matrix}

\[
K^\top=
\begin{bmatrix}
1&0&1\\
1&1&2
\end{bmatrix},
\]
so
\[
S=QK^\top=
\begin{bmatrix}
1&0&1\\
1&1&2\\
2&1&3
\end{bmatrix}.
\]

Scaling by $\sqrt{d_k}=\sqrt{2}$ gives
\[
S_{\mathrm{scaled}}
=\frac{S}{\sqrt{2}}
\approx
\begin{bmatrix}
0.7071&0&0.7071\\
0.7071&0.7071&1.4142\\
1.4142&0.7071&2.1213
\end{bmatrix}.
\]
The score matrices have shape $3\times3$.

\section*{3. Apply row-wise softmax}

For a row $z$, softmax is
\[
\softmax(z)_j=\frac{e^{z_j}}{\sum_\ell e^{z_\ell}}.
\]
Applying it independently to each row gives
\[
A\approx
\begin{bmatrix}
0.4011&0.1978&0.4011\\
0.2483&0.2483&0.5035\\
0.2840&0.1400&0.5760
\end{bmatrix}.
\]
Each row sums to approximately $1$; tiny deviations are only rounding error.

For example, the `cats' row is
\[
\begin{bmatrix}0.2840&0.1400&0.5760\end{bmatrix},
\]
so `cats' assigns about $28.4\%$ of its attention to `I', $14.0\%$ to `like', and $57.6\%$ to itself.

\section*{4. Compute the output}

\[
O=AV
\approx
\begin{bmatrix}
1.4011&0.5989\\
1.5035&0.7517\\
1.5760&0.7160
\end{bmatrix}.
\]
The output has shape $3\times2$.

For the final row,
\[
\begin{aligned}
o_{\mathrm{cats}}
&=0.2840\begin{bmatrix}1&0\end{bmatrix}
+0.1400\begin{bmatrix}1&1\end{bmatrix}
+0.5760\begin{bmatrix}2&1\end{bmatrix}\\
&\approx\begin{bmatrix}1.5760&0.7160\end{bmatrix}.
\end{aligned}
\]

\section*{Final answer}

For `cats', the largest attention weight is the self-attention entry on `cats' itself:
\[
\boxed{0.5760}.
\]

\section*{Bonus}

No general relationship is required between the largest attention weight and the value vector with the largest norm. The attention weights are determined by $Q$ and $K$ through
\[
A=\softmax\left(\frac{QK^\top}{\sqrt{d_k}}\right),
\]
while $V$ is used only afterward in
\[
O=AV.
\]
A token can therefore receive the largest attention weight even when its value vector does not have the largest norm, and vice versa. Any coincidence in a particular toy example is not a structural requirement of attention.

\section*{Pipeline summary}

\[
X\longrightarrow(Q,K,V)
\longrightarrow QK^\top
\longrightarrow \frac{QK^\top}{\sqrt{d_k}}
\longrightarrow \softmax
\longrightarrow A
\longrightarrow AV.
\]

\section*{Numerical verification}

The companion \texttt{answer\_key.py} reproduces these matrices with PyTorch and asserts the expected values, shapes, and row sums.

\end{document}
